\documentclass[12pt,a4paper]{article}
\usepackage[margin=2.5cm]{geometry}
\usepackage{amsmath,amssymb}
\usepackage{booktabs}
\usepackage{xcolor}
\usepackage{hyperref}
\usepackage{enumitem}
\usepackage{titlesec}
\usepackage{mdframed}
\usepackage{parskip}
\usepackage{graphicx}
\usepackage{fancyhdr}
\usepackage{lmodern}
\usepackage[T1]{fontenc}

\hypersetup{colorlinks=true, linkcolor=blue!70!black, urlcolor=blue}

% ---- Colours ----
\definecolor{headerblue}{RGB}{0,70,127}
\definecolor{accentgreen}{RGB}{0,120,80}
\definecolor{warnred}{RGB}{180,30,30}
\definecolor{lightgray}{RGB}{245,245,245}
\definecolor{boxblue}{RGB}{230,240,255}

% ---- Section styling ----
\titleformat{\section}{\large\bfseries\color{headerblue}}{}{0em}{}[\titlerule]
\titleformat{\subsection}{\normalsize\bfseries\color{headerblue!80!black}}{}{0em}{}

% ---- Header / footer ----
\pagestyle{fancy}
\fancyhf{}
\renewcommand{\headrulewidth}{0.4pt}
\fancyhead[L]{\small\color{gray} Research Defence Guide --- Kusal Pabasara}
\fancyhead[R]{\small\color{gray} ICDSIAI-26 Accepted Paper}
\fancyfoot[C]{\small\thepage}

% ---- Key number box ----
\newmdenv[backgroundcolor=boxblue, linecolor=headerblue, linewidth=1.2pt,
          innertopmargin=8pt, innerbottommargin=8pt,
          innerleftmargin=10pt, innerrightmargin=10pt,
          skipabove=6pt, skipbelow=6pt]{keybox}

% ---- Q&A box ----
\newmdenv[backgroundcolor=lightgray, linecolor=accentgreen, linewidth=1.2pt,
          innertopmargin=8pt, innerbottommargin=8pt,
          innerleftmargin=10pt, innerrightmargin=10pt,
          skipabove=6pt, skipbelow=6pt]{qabox}

% ---- Warning box ----
\newmdenv[backgroundcolor=red!5, linecolor=warnred, linewidth=1.2pt,
          innertopmargin=8pt, innerbottommargin=8pt,
          innerleftmargin=10pt, innerrightmargin=10pt,
          skipabove=6pt, skipbelow=6pt]{warnbox}

\begin{document}

% ============================================================
%  TITLE PAGE
% ============================================================
\begin{titlepage}
\centering
\vspace*{2cm}
{\Huge\bfseries\color{headerblue} Research Defence Guide}\\[0.5cm]
{\large\color{gray} Everything you need to know to confidently present,\\
defend, and secure sponsorship for your research}\\[1.5cm]
\rule{0.8\textwidth}{1pt}\\[0.8cm]
{\Large\bfseries Interest-Driven Personalized Sports Recommendation\\
Using Classical Ensemble Learning with a\\
Novel-Sport Discovery Mechanism}\\[0.8cm]
\rule{0.8\textwidth}{0.5pt}\\[1cm]
{\large Wandawa Gamage Kusal Pabasara}\\[0.2cm]
{\normalsize Department of Computer Science and Engineering\\
University of Moratuwa, Sri Lanka\\
\texttt{kusalp.23@cse.mrt.ac.lk}}\\[1cm]
{\normalsize\color{accentgreen}\bfseries Accepted at ICDSIAI-26}\\[0.3cm]
{\normalsize International Conference on Data Science Integration\\
with Artificial Intelligence}\\[2cm]
{\small\color{gray} Prepared: March 2026 \quad|\quad
GitHub: \url{https://github.com/KusalPabasara/sports-recommendation}}
\end{titlepage}

% ============================================================
%  TABLE OF CONTENTS
% ============================================================
\tableofcontents
\newpage

% ============================================================
\section{The Big Picture --- What Is This Research About?}
% ============================================================

\subsection{The core idea in one sentence}

\begin{keybox}
\textbf{This research builds a machine learning system that recommends
sports to people based on their \emph{personal interests and personality},
not just their physical fitness --- and it also suggests sports they have
never tried before.}
\end{keybox}

\subsection{The real-world problem}

Imagine you are a school student or a general adult who wants to find a sport
that suits you. The technology that exists today works like this: it measures
how fast you run, how high you jump, and how flexible you are, then tells you
which sport you are \emph{physically suited for}. Apps like AI Sport Scout do
exactly this.

\textbf{The problem is that physical fitness is not the main reason people
stick with a sport.} Sports psychology has shown clearly and repeatedly that
the reason people keep playing a sport long-term is because they \emph{enjoy
it} --- it matches their personality, their interests, and their social
preferences. If someone hates team activities, telling them they are
physically suited for football will not make them enjoy football.

The WHO (World Health Organization) reported in 2022 that over \textbf{1.4 billion
adults worldwide are insufficiently physically active}. One major reason is
that people do not find sports that genuinely suit them. This research
directly addresses that gap.

\subsection{The research question}

\begin{keybox}
\textit{``Can a machine learning system that prioritizes personal interests
over physical traits provide better and more holistic sports recommendations
to general users?''}
\end{keybox}

The answer this research found: \textbf{Yes, and it outperforms physical-only
systems by 17.4\% on the standard ranking quality metric (NDCG@5).}

% ============================================================
\section{Five Key Contributions --- What Is New Here?}
% ============================================================

Your paper makes five original contributions. Know these by heart.

\begin{enumerate}[leftmargin=*, label=\textbf{\arabic*.}]

\item \textbf{First interest-first sports recommendation framework.}\\
Previous systems all put physical metrics first. This is the first system
to put personal interest scores (Likert-scale survey answers) as the
\emph{primary} input signal, with physical traits as secondary. This is a
fundamental design philosophy shift.

\item \textbf{Stacking ensemble with 17.4\% improvement.}\\
A stacking ensemble combining XGBoost + Random Forest (base learners) with
Logistic Regression (meta-learner) outperforms all individual models and all
physical-only baselines by 17.4\% on NDCG@5. Only classical machine learning
methods are used --- no deep learning, no neural networks.

\item \textbf{Novel discovery score for unexplored sports.}\\
A cosine-similarity-based discovery score that identifies sports a user has
\emph{never tried} but is likely to enjoy based on latent interest alignment.
This achieves an 85.7\% novelty rate --- meaning 85.7\% of its top-5
recommendations are sports the user has not encountered before.

\item \textbf{SHAP validation of interest-first design.}\\
Using SHAP (a model explanation tool), it was proven that interest features
account for 65.9\% of the model's predictive power when combined with
self-rated strength features. This is not just a design choice --- it is
empirically validated.

\item \textbf{Public dataset and full code release.}\\
A literature-calibrated synthetic dataset of 1,000 users across 20 sports,
with all correlations verified against published psychology research, is
released publicly on GitHub for other researchers to use.

\end{enumerate}

% ============================================================
\section{Literature Review --- What Others Have Done and Why It Is Not Enough}
% ============================================================

You reviewed five areas of prior work. Here is what to say about each:

\subsection{Sports talent identification (e.g., AI Sport Scout, Ramos et al.)}

Systems like AI Sport Scout use battery fitness tests (sprint times, jump
distances, flexibility) to match athletes to sports. While effective for
\emph{elite scouting}, they:
\begin{itemize}
\item Require specialized physical testing equipment --- inaccessible to
general users
\item Assume physical aptitude is the primary signal
\item Do not model interest, motivation, or personality at all
\item Have a ceiling: physical features alone explain only $\sim$60\% of
variance in sports participation (Lidor et al., 2005)
\end{itemize}

\subsection{Recommendation systems (Burke, Zhang et al.)}

\begin{sloppypar}
General recommendation systems (collaborative filtering, content-based) are
well-established but suffer from a \textbf{cold-start problem} for novice
users who have no prior sports history. Deep learning recommenders improve
expressiveness but do not incorporate physical or interest features.
\end{sloppypar}

\subsection{Sports psychology (Vallerand, Pelletier, Fraser-Thomas)}

This is the theoretical backbone of the research. Three landmark findings:
\begin{itemize}
\item \textbf{Vallerand et al. (2003)}: Harmonious passion (interest-driven)
predicts sustained engagement; obsessive passion leads to burnout.
\item \textbf{Pelletier et al. (1995)}: Intrinsic motivation driven by
interest is the strongest longitudinal predictor of sport continuation.
\item \textbf{Fraser-Thomas \& C\^{o}t\'{e} (2006)}: Interest and enjoyment
outperform physical talent in predicting youth sport persistence. Mismatch
between assigned sports and personal interest is the primary driver of dropout.
\end{itemize}
\textbf{These three studies are the scientific justification for why this
research puts interests first.}

\subsection{Student/youth sports ML (IARCI, Holt)}

The closest prior work on ML for student sports (IARCI, 2023) used decision
trees and Na\"{i}ve Bayes on small samples. It classifies a single preferred
sport from 3--5 options. It does not do multi-label recommendation across 20
sports and has no discovery mechanism. This research fills that gap.

\subsection{Ensemble learning (Wolpert 1992, Breiman 2001)}

Stacking ensembles were introduced by Wolpert (1992) and have been proven
effective when base learners have complementary strengths. XGBoost handles
non-linear interactions well; Random Forest reduces variance through bagging.
Combining them through a logistic regression meta-learner is theoretically
sound and practically strong.

% ============================================================
\section{Methodology --- Every Component Explained Simply}
% ============================================================

\subsection{What does the model take as input?}

Each user is described by \textbf{29 features} in four groups:

\begin{center}
\begin{tabular}{llp{7cm}}
\toprule
\textbf{Group} & \textbf{Count} & \textbf{Examples} \\
\midrule
Interests (Likert 1--5) & 12 & Team vs.\ individual, outdoor/indoor, competition drive \\
Strengths (Likert 1--5) & 8 & Self-rated endurance, speed, flexibility, coordination \\
Physical metrics & 6 & Height, weight, BMI, sprint time, jump distance, age \\
Demographics & 3 & Gender, region, facility access \\
\midrule
\textbf{Total} & \textbf{29} & \\
\bottomrule
\end{tabular}
\end{center}

\subsection{What does the model predict?}

For each of 20 sports, the model predicts:
\begin{itemize}
\item \textbf{Play probability}: Will this user actively play this sport?
\item \textbf{Watch probability}: Will this user watch/follow this sport?
\end{itemize}
This is a \emph{multi-label binary classification} problem (20 binary
predictions per user per task).

\subsection{The dataset --- why synthetic is acceptable}

No public dataset exists that combines interest profiles, physical metrics,
and multi-label sports engagement labels for general users. Therefore, a
synthetic dataset was generated following three principles:

\begin{enumerate}
\item \textbf{Literature-calibrated}: Correlations between features and
labels are set to match published benchmarks:
  \begin{itemize}
  \item Interest--enjoyment Pearson $r \geq 0.60$ (Vallerand et al., 2003)
  \item Interest--hours/week $r \geq 0.50$ (Pelletier et al., 1995)
  \item Extraversion positively correlated with team sport participation
        (Allen et al., 2013)
  \end{itemize}
\item \textbf{Reproducible}: All generation code is public on GitHub.
\item \textbf{Precedented}: Using synthetic data for feasibility studies in
data-scarce domains is an accepted practice, explicitly stated in the paper's
limitations section.
\end{enumerate}

The dataset has 1,000 users, 20 sports, and is split 80/20 (800 train,
200 test) with stratified sampling.

\subsection{The affinity score --- how labels are generated}

The raw likelihood of a user $u$ engaging with sport $s$ is:
\[
a_{u,s} = \underbrace{2 \cdot \frac{\mathbf{i}_u \cdot \mathbf{p}_s}{D_I}}_{\text{interest term (weight 2)}}
         + \underbrace{0.5 \cdot \phi_{u,s}}_{\text{physical term (weight 0.5)}}
         + b_s + \epsilon
\]
The factor of 2 on the interest term and 0.5 on the physical term
\textbf{operationalizes the interest-first weighting} directly in data
generation, ensuring the synthetic data reflects the design philosophy.

\subsection{Preprocessing}

\begin{itemize}
\item Continuous features (age, height, weight, BMI, sprint, jump):
      \textbf{z-score normalized} on training data only (no data leakage)
\item Categorical features (gender, region): integer-encoded
\item Likert features: kept as ordinal integers (1--5 scale preserved)
\end{itemize}

\subsection{Hyperparameter tuning}

Grid search with 3-fold stratified cross-validation was applied to both
XGBoost and Random Forest. The search spaces were:

\begin{center}
\begin{tabular}{lll}
\toprule
\textbf{Model} & \textbf{Parameter} & \textbf{Search Space} \\
\midrule
XGBoost & n\_estimators & \{100, 200, 300\} \\
        & max\_depth & \{3, 5, 7\} \\
        & learning\_rate & \{0.01, 0.05, 0.1\} \\
        & subsample & \{0.7, 0.8, 1.0\} \\
        & colsample\_bytree & \{0.7, 0.8, 1.0\} \\
\midrule
Random Forest & n\_estimators & \{100, 200, 300\} \\
              & max\_depth & \{5, 8, 12, None\} \\
              & min\_samples\_leaf & \{2, 5, 10\} \\
              & max\_features & \{$\sqrt{d}$, $\log_2 d$, $0.5d$\} \\
\bottomrule
\end{tabular}
\end{center}

\textbf{Best parameters used:} XGBoost (n\_estimators=200, max\_depth=5,
lr=0.05); Random Forest (n\_estimators=200, max\_depth=8).

\subsection{The stacking ensemble}

\begin{keybox}
\textbf{Architecture:}
\begin{enumerate}
\item XGBoost and Random Forest each make predictions on out-of-fold data
(3-fold CV), generating probability estimates
\item These probability estimates are concatenated and fed as input features
to a Logistic Regression meta-learner
\item The meta-learner learns to optimally weight and combine the two base
learners' predictions
\end{enumerate}
\textbf{Why this works:} XGBoost captures complex non-linear interactions
between features. Random Forest reduces variance through averaging many
decorrelated trees. Logistic Regression provides a simple, interpretable
combination that does not overfit. Each model contributes what it is
best at.
\end{keybox}

The entire pipeline is wrapped in \texttt{MultiOutputClassifier} to handle
all 20 sports simultaneously.

\subsection{The discovery score}

The discovery score answers: \textit{``Which sports has this user never tried
but would likely enjoy based on their interests?''}

\textbf{Step 1 --- Build a sport profile:} For each sport $s$, compute the
average interest vector of all users in the training set who play that sport:
\[
\mathbf{q}_s = \frac{1}{|\mathcal{U}_s|} \sum_{u \in \mathcal{U}_s} \mathbf{i}_u
\]
This creates a ``fingerprint'' of what kind of person tends to enjoy each sport.

\textbf{Step 2 --- Score unexplored sports:} For a new user, compute cosine
similarity between their interest vector and each sport's fingerprint, but
only for sports they have \emph{not tried before}:
\[
d_{u,s} = \frac{\mathbf{i}_u \cdot \mathbf{q}_s}{\|\mathbf{i}_u\|\,\|\mathbf{q}_s\|}
           \cdot (1 - \text{tried}_{u,s})
\]

\textbf{Step 3 --- Blend with model prediction:}
\[
\hat{r}_{u,s} = \alpha \cdot \tilde{d}_{u,s} + (1-\alpha) \cdot \tilde{f}_{u,s}
\]
where $\alpha = 0.4$ (40\% discovery, 60\% model prediction). Setting
$\alpha \to 0$ gives pure model recommendations; $\alpha \to 1$ maximizes
novelty. This is a controllable dial.

% ============================================================
\section{Baselines --- What You Are Compared Against}
% ============================================================

Six baselines were implemented to rigorously test the system:

\begin{center}
\begin{tabular}{llp{8.5cm}}
\toprule
\textbf{ID} & \textbf{Name} & \textbf{What it represents} \\
\midrule
B1 & Random & Lower bound --- completely random \\
B2 & Popularity & Recommends most popular sports to everyone --- zero personalization \\
B3 & Physical-only & XGBoost on 6 physical features only (AI Sport Scout approach) \\
B4 & Interest-only & XGBoost on 12 interest features only --- ablates physical features \\
B5 & Full XGBoost & XGBoost on all 29 features, no tuning, no discovery \\
B6 & Full RF & Random Forest on all 29 features, no tuning, no discovery \\
\midrule
T1 & Tuned XGBoost & Hyperparameter-tuned XGBoost \\
T2 & Tuned RF & Hyperparameter-tuned Random Forest \\
\textbf{S1} & \textbf{Stacking} & \textbf{Our main model --- stacking ensemble} \\
S2 & Stack+Discovery & Stacking ensemble + discovery module activated \\
\bottomrule
\end{tabular}
\end{center}

% ============================================================
\section{Results --- Every Number and What It Means}
% ============================================================

\subsection{Evaluation metrics explained}

\begin{itemize}
\item \textbf{NDCG@5} (Normalized Discounted Cumulative Gain): Measures
ranking quality --- are the truly relevant sports ranked higher in the top 5?
Higher is better. Range 0--1. \textbf{This is the primary metric.}
\item \textbf{MAP@5} (Mean Average Precision): Measures precision across
ranking positions. Higher is better.
\item \textbf{P@5} (Precision@5): What fraction of the top-5 recommendations
are actually relevant?
\item \textbf{R@5} (Recall@5): What fraction of all relevant sports appear
in the top 5?
\item \textbf{DR@5} (Discovery Rate@5): What fraction of top-5 are sports
the user has never tried? Only S2 optimizes this.
\end{itemize}

\subsection{Play task results (main table)}

\begin{center}
\begin{tabular}{lrrrrr}
\toprule
\textbf{Model} & \textbf{P@5} & \textbf{R@5} & \textbf{NDCG@5} & \textbf{MAP@5} & \textbf{DR@5} \\
\midrule
B1 Random         & 0.201 & 0.247 & 0.228 & 0.500 & 0.799 \\
B2 Popularity     & 0.380 & 0.487 & 0.502 & 0.715 & 0.620 \\
B3 Physical-only  & 0.335 & 0.415 & 0.424 & 0.685 & 0.665 \\
B4 Interest-only  & 0.321 & 0.419 & 0.406 & 0.629 & 0.679 \\
B5 Full XGBoost   & 0.332 & 0.422 & 0.429 & 0.678 & 0.668 \\
B6 Full RF        & 0.372 & 0.478 & 0.491 & 0.708 & 0.628 \\
\midrule
T1 Tuned XGBoost  & 0.317 & 0.401 & 0.413 & 0.677 & 0.683 \\
T2 Tuned RF       & 0.371 & 0.476 & 0.493 & \textbf{0.725} & 0.629 \\
\textbf{S1 Stacking} & \textbf{0.381} & \textbf{0.487} & \textbf{0.498} & 0.705 & 0.619 \\
S2 Stack+Discovery & 0.143 & 0.179 & 0.203 & 0.639 & \textbf{0.857} \\
\bottomrule
\end{tabular}
\end{center}

\textbf{Key numbers to remember:}

\begin{keybox}
\begin{itemize}
\item S1 NDCG@5 = \textbf{0.498} vs B3 Physical-only = 0.424
      $\Rightarrow$ \textbf{17.4\% improvement} over the current paradigm
\item S1 NDCG@5 = 0.498 vs B5 Full XGBoost = 0.429
      $\Rightarrow$ \textbf{16.1\% improvement} over untuned full model
\item S2 DR@5 = \textbf{0.857} --- 85.7\% of top-5 are sports never tried before
\item T2 achieves highest MAP@5 = \textbf{0.725} ---
      hyperparameter tuning is worth it for ranking precision
\end{itemize}
\end{keybox}

\subsection{Why B2 Popularity scores high and what that means}

B2 Popularity (NDCG@5 = 0.502) slightly beats S1 in NDCG. \textbf{This is
not a weakness --- it is expected and explained.} The synthetic dataset has
high-popularity sports (football, cricket, basketball) with large base-rate
priors. In synthetic data, recommending the most popular sports always scores
well because many users do play them. In \emph{real-world} data with greater
user diversity and niche preferences, personalized models consistently
outperform popularity baselines. This is explicitly discussed in the paper.

\subsection{Watch task results}

\begin{center}
\begin{tabular}{lrrrr}
\toprule
\textbf{Model} & \textbf{P@5} & \textbf{R@5} & \textbf{NDCG@5} & \textbf{MAP@5} \\
\midrule
B1 Random         & 0.555 & 0.267 & 0.561 & 0.716 \\
B2 Popularity     & 0.724 & 0.352 & 0.741 & 0.851 \\
B5 Full XGBoost   & 0.656 & 0.315 & 0.677 & 0.814 \\
\textbf{S1 Stacking} & \textbf{0.693} & \textbf{0.335} & \textbf{0.721} & \textbf{0.852} \\
\bottomrule
\end{tabular}
\end{center}

The stacking ensemble outperforms XGBoost on the watch task as well,
confirming the ensemble \textbf{generalizes across engagement levels}
(playing vs.\ watching), not just on the primary play task.

\subsection{Ablation study --- proving every component matters}

The ablation study systematically removes components to show each contributes:

\begin{center}
\begin{tabular}{lrrr}
\toprule
\textbf{Configuration} & \textbf{NDCG@5} & \textbf{MAP@5} & \textbf{DR@5} \\
\midrule
Physical-only           & 0.424 & 0.685 & 0.665 \\
Interest-only           & 0.406 & 0.629 & 0.679 \\
Full XGBoost (all feat) & 0.429 & 0.678 & 0.668 \\
Full RF (all feat)      & 0.491 & 0.708 & 0.628 \\
Tuned XGBoost           & 0.413 & 0.677 & 0.683 \\
Tuned RF                & 0.493 & 0.725 & 0.629 \\
\textbf{Stacking (S1)}  & \textbf{0.498} & 0.705 & 0.619 \\
S1 + Discovery (S2)     & 0.203 & 0.639 & \textbf{0.857} \\
\bottomrule
\end{tabular}
\end{center}

\textbf{What this proves:}
\begin{enumerate}
\item \textbf{Features matter}: Both interest-only (0.406) and physical-only
(0.424) are weaker than the full model (0.498). Combining all features helps.
\item \textbf{Stacking beats individuals}: S1 (0.498) outperforms both T1
(0.413) and T2 (0.493) individually --- the ensemble is genuinely better.
\item \textbf{Discovery trades accuracy for novelty by design}: S2 (0.203
NDCG, 0.857 DR) --- the low accuracy is intentional, as S2 targets unexplored
sports, not the most certain predictions.
\end{enumerate}

\subsection{SHAP feature importance}

SHAP (SHapley Additive exPlanations) is a mathematically principled tool
from game theory that quantifies how much each feature contributes to each
prediction. Results averaged over 200 test users, three sports:

\begin{center}
\begin{tabular}{lrrrr}
\toprule
\textbf{Feature Group} & \textbf{Football} & \textbf{Swimming} & \textbf{Esports} & \textbf{Mean} \\
\midrule
Interests          & 41.4\% & 30.5\% & 34.2\% & \textbf{35.4\%} \\
Strengths (self)   & 27.8\% & 40.6\% & 23.2\% & \textbf{30.5\%} \\
Physical (objective) & 26.2\% & 23.4\% & 30.2\% & \textbf{26.6\%} \\
Demographics       & 4.6\%  & 5.6\%  & 12.4\% & \textbf{7.5\%} \\
\midrule
\textbf{Interest + Strength} & \textbf{69.2\%} & \textbf{71.1\%} & \textbf{57.4\%} & \textbf{65.9\%} \\
\bottomrule
\end{tabular}
\end{center}

\begin{keybox}
Interest + Strength (psychological signal) = \textbf{65.9\%} of total
predictive power on average. Physical metrics = only 26.6\%.\\[4pt]
\textbf{This empirically proves the interest-first design is correct.}
The model itself ``learned'' that interests matter more --- this was not
forced; it emerged from the data.
\end{keybox}

% ============================================================
\section{Limitations --- Be Honest and Have Answers Ready}
% ============================================================

\begin{warnbox}
Always acknowledge limitations \emph{before} being asked --- it shows
academic maturity. Then immediately explain the response.
\end{warnbox}

\begin{enumerate}[leftmargin=*, label=\textbf{L\arabic*.}]

\item \textbf{Synthetic data.}\\
\textit{Limitation}: The dataset is computer-generated, not from a real survey.\\
\textit{Response}: ``Synthetic data is an accepted and common approach in
feasibility studies, especially in data-scarce domains where no public dataset
exists. All correlations in the synthetic data are calibrated to match actual
published studies in sports psychology. The paper explicitly frames results as
a proof-of-concept and identifies a real-user survey (N$\geq$500) as the
critical next step. The framework is technically validated; the next phase is
empirical validation.''

\item \textbf{Popularity baseline nearly matches our model.}\\
\textit{Limitation}: B2 Popularity (0.502) slightly outperforms S1 (0.498) in NDCG.\\
\textit{Response}: ``This is a known and expected property of synthetic data.
High-popularity sports have inflated base rates. In real-world deployments
with heterogeneous user populations and niche interests, personalized models
consistently outperform popularity baselines. The paper explains this directly
in the Discussion section.''

\item \textbf{20-sport taxonomy only.}\\
\textit{Limitation}: Only 20 sports are covered.\\
\textit{Response}: ``The 20 sports cover the major global categories: team,
individual, indoor, outdoor, high-intensity, and low-intensity. The framework
architecture scales trivially to any number of sports --- adding more sports
simply adds more output labels to the MultiOutputClassifier. The 20-sport
scope was chosen for tractability in this proof-of-concept stage.''

\item \textbf{Cold-start for entirely new sports.}\\
\textit{Limitation}: If a sport has no training users, the discovery score
falls back to the training set mean.\\
\textit{Response}: ``This is a standard cold-start limitation in
content-collaborative hybrid systems, well documented in the recommender
systems literature. Hand-crafted sport profile vectors (like those used in
affinity score generation) can be used as fallback for new sports.''

\item \textbf{Self-report bias in interest features.}\\
\textit{Limitation}: Interest and strength features are self-rated on Likert
scales; people may mis-rate themselves.\\
\textit{Response}: ``Self-report Likert scales are the standard instrument in
sports psychology research (Vallerand, Pelletier). Their validity is
well-established in the literature this work is grounded in. Cross-validation
against behavioral data (e.g., exercise logs) is identified as future work.''

\end{enumerate}

% ============================================================
\section{Tough Questions Your Professor Will Ask --- With Answers}
% ============================================================

\subsection{Questions about methodology}

\begin{qabox}
\textbf{Q: Why did you use classical ML and not deep learning?}\\[4pt]
\textbf{A:} ``Three reasons. First, the dataset has 1,000 users and 29
features --- deep learning typically needs far more data to outperform
classical methods. Second, classical methods like XGBoost and Random Forest
are inherently interpretable and allow SHAP analysis, which is essential for
validating the interest-first hypothesis. Third, for a deployment setting
(schools, fitness apps), interpretable models are preferred for trust and
regulatory compliance. A system that recommends sports needs to be explainable
to users.''
\end{qabox}

\begin{qabox}
\textbf{Q: Why is the dataset synthetic? Isn't that a weakness?}\\[4pt]
\textbf{A:} ``It is a limitation, not a disqualifying weakness. No public
dataset combines multi-label sports engagement with interest profiles and
physical metrics. Creating a synthetic dataset calibrated to published
psychology findings is an accepted approach for feasibility studies in
data-scarce domains. The paper is transparent about this, frames the work
as a proof-of-concept, and identifies real-user data collection as the
primary next step. The conference reviewers accepted this framing.''
\end{qabox}

\begin{qabox}
\textbf{Q: Why a stacking ensemble specifically? Why not just use XGBoost?}\\[4pt]
\textbf{A:} ``XGBoost alone achieves NDCG@5 = 0.429 (B5). The stacking
ensemble achieves 0.498 --- a 16.1\% improvement. The ablation study
isolates this: every component --- feature combination, tuning, stacking ---
contributes to the final result. XGBoost and Random Forest have complementary
strengths: XGBoost captures non-linear feature interactions while Random Forest
reduces variance through bagging. The meta-learner learns when to trust each.
The results empirically justify the choice.''
\end{qabox}

\begin{qabox}
\textbf{Q: What is NDCG and why is it the right metric?}\\[4pt]
\textbf{A:} ``NDCG (Normalized Discounted Cumulative Gain) measures ranking
quality. In a recommendation system, it is not enough to just return the right
items --- you need to rank them correctly, with the most relevant items at the
top. NDCG gives more credit for relevant items appearing at position 1 than at
position 5. It is the standard primary metric for ranking-based recommendation
systems, used in virtually all recommendation papers.''
\end{qabox}

\begin{qabox}
\textbf{Q: Why is $\alpha = 0.4$ chosen for the discovery blend?}\\[4pt]
\textbf{A:} ``Alpha of 0.4 means 40\% weight on discovery and 60\% on model
prediction. This was chosen as a balanced default that gives meaningful novelty
without completely sacrificing recommendation accuracy. It is a design
parameter, not a tuned hyperparameter. The paper explicitly states that alpha
can be adapted per user based on their stated exploration preference. Future
work includes learning alpha from user feedback.''
\end{qabox}

\begin{qabox}
\textbf{Q: Why does Tuned XGBoost underperform untuned XGBoost?}\\[4pt]
\textbf{A:} ``This is discussed openly in the paper. The grid search was
performed on the first sport label (highest base rate) and applied to all 20
sports. This proxy-label tuning strategy can overfit the hyperparameters to
the dominant sport and not generalize well across all 20. Per-sport tuning is
identified as a direction for future work. The fact that the stacking ensemble
(S1 = 0.498) still outperforms both demonstrates that the ensemble approach
is robust even when individual components are not perfectly tuned.''
\end{qabox}

\subsection{Questions about the research contribution}

\begin{qabox}
\textbf{Q: What is actually new here? Recommender systems exist already.}\\[4pt]
\textbf{A:} ``The novelty is in the combination of four things no prior
system has done simultaneously: (1) interest-first design philosophy where
psychological interest features are the primary input, not physical metrics;
(2) accessibility for novice and general users who have no sports history;
(3) a discovery mechanism that actively surfaces unexplored sports; and
(4) using only classical ML methods with full explainability. The Research Gap
table in the paper shows this gap clearly against five prior systems.''
\end{qabox}

\begin{qabox}
\textbf{Q: The SHAP results show 65.9\% for interests+strengths. But strengths
are also somewhat physical. Is this really showing interests are important?}\\[4pt]
\textbf{A:} ``That is a fair point and is addressed in the paper. Strength
features are Likert-scaled self-ratings of perceived physical capability ---
not objective measurements like height or sprint time. They are psychological
assessments of physical ability, adjacent to how a person sees themselves
rather than what a clinician measures. This makes them \emph{psychologically}
adjacent features, distinct from the six objective physical metrics. The paper
is explicit about this framing.''
\end{qabox}

\begin{qabox}
\textbf{Q: Has this been tested on real users?}\\[4pt]
\textbf{A:} ``Not yet --- and the paper is completely transparent about this.
The current work is a proof-of-concept demonstrating that the framework is
technically sound and internally consistent with published psychology findings.
A randomized A/B user study (N=60, 30 per arm) is described in detail in the
paper as the planned immediate next step, with a pre-specified primary outcome
(sports engagement intention score at two-week follow-up) and statistical
power calculations ($\beta$ = 0.80, $p < 0.05$).''
\end{qabox}

\subsection{Questions about the conference and publication}

\begin{qabox}
\textbf{Q: Why ICDSIAI-26 specifically?}\\[4pt]
\textbf{A:} ``ICDSIAI-26 focuses on the intersection of data science and
artificial intelligence, which is exactly where this work sits --- it is an
ML system grounded in psychology and applied to a real-world public health
problem. The conference audience includes researchers in applied ML,
recommender systems, and human-centered AI, which are the communities most
relevant to this work. The paper has been accepted after peer review.''
\end{qabox}

\begin{qabox}
\begin{sloppypar}
\textbf{Q: How do you know this work is original and not plagiarized?}\\[4pt]
\textbf{A:} ``The code, dataset, and paper are all original work. The full
implementation history is on GitHub at
\url{https://github.com/KusalPabasara/sports-recommendation}.
The paper was submitted through the conference's plagiarism-checking process
and was accepted. All prior work is properly cited with 21 references across
sports analytics, recommendation systems, psychology, and ensemble learning.''
\end{sloppypar}
\end{qabox}

% ============================================================
\section{Sponsorship Case --- How to Argue for Funding}
% ============================================================

\subsection{What you are asking for}

Conference registration fee + travel/accommodation to present the accepted
paper at ICDSIAI-26.

\subsection{Arguments to make to your professor / department}

\begin{enumerate}[leftmargin=*, label=\textbf{Arg \arabic{enumi}.}]

\item \textbf{Peer-reviewed acceptance --- the work is externally validated.}\\
``This paper has been peer-reviewed and accepted by ICDSIAI-26. Acceptance
means independent expert reviewers found the research novel, technically
sound, and of sufficient quality for publication. This is not a self-assessed
claim --- it is external academic validation.''

\item \textbf{Institutional visibility.}\\
``The affiliation on the paper is the University of Moratuwa, Department of
Computer Science and Engineering. Every publication citing this work will
cite the university. This is direct, measurable research output and
international visibility for the department.''

\item \textbf{A rare undergraduate achievement.}\\
``First-authored peer-reviewed international conference papers at
undergraduate level are rare. Sponsoring this presentation demonstrates the
department's commitment to undergraduate research and will strengthen the
department's research profile.''

\item \textbf{Real-world impact potential.}\\
``Sports participation is a global public health problem. WHO reports 1.4
billion insufficiently active adults. This research addresses a specific,
unmet need: helping novice and general users discover sports that genuinely
suit them. This has real deployment potential for schools, national sports
bodies, fitness apps, and public health programs.''

\item \textbf{Follow-on research funding potential.}\\
``Presenting at an international conference creates connections with other
researchers and potential collaborators. The planned user study (N=60 A/B
randomized trial) described in the paper would be a strong candidate for
a follow-on research grant, which would bring external funding to the
department.''

\item \textbf{Dataset and code are publicly released.}\\
``The synthetic dataset and full code are already public on GitHub. This is
research that other groups can build on --- maximizing the research impact
beyond the single publication.''

\end{enumerate}

\subsection{Suggested email / talking points structure}

\begin{mdframed}[backgroundcolor=lightgray]
\textbf{To Professor [Name]:}

I am writing to request sponsorship for conference participation at
ICDSIAI-26, where my paper titled \emph{``Interest-Driven Personalized
Sports Recommendation Using Classical Ensemble Learning with a Novel-Sport
Discovery Mechanism''} has been accepted for publication.

The paper presents a machine learning system that recommends sports to users
based on personal interests rather than physical fitness, achieving a 17.4\%
improvement over the current physical-only paradigm. It was conducted as
independent research and has been peer-reviewed and accepted at ICDSIAI-26.

The paper is affiliated with the Department of Computer Science and
Engineering, University of Moratuwa, providing direct international
visibility for the department. I would be grateful for sponsorship
covering the conference registration [and travel/accommodation if applicable].

I am happy to present the full research to you at your convenience before
the conference to discuss the work in detail.

\medskip
Yours sincerely,\\
Kusal Pabasara
\end{mdframed}

% ============================================================
\section{Quick Reference Card --- Numbers to Know Instantly}
% ============================================================

\begin{keybox}
\textbf{Memorize these. They will come up in any conversation about the research.}

\medskip
\begin{tabular}{lp{9cm}}
\textbf{Dataset} & 1,000 users, 20 sports, 29 features, 800 train / 200 test \\
\textbf{Best model} & S1 Stacking (XGBoost + RF + Logistic Regression meta-learner) \\
\textbf{Best NDCG@5} & 0.498 (S1 Stacking) \\
\textbf{vs Physical-only} & 17.4\% improvement (0.424 $\to$ 0.498) \\
\textbf{vs Full XGBoost} & 16.1\% improvement (0.429 $\to$ 0.498) \\
\textbf{Best MAP@5} & 0.725 (T2 Tuned Random Forest) \\
\textbf{Discovery rate} & 85.7\% novelty at K=5 (S2 Stack+Discovery) \\
\textbf{SHAP interest} & 35.4\% interests alone; 65.9\% interests + strengths \\
\textbf{SHAP physical} & 26.6\% objective physical metrics \\
\textbf{Watch task S1} & NDCG@5 = 0.721, MAP@5 = 0.852 \\
\textbf{Discovery blend} & $\alpha = 0.4$ (40\% discovery, 60\% model) \\
\textbf{Grid search} & 3-fold stratified CV \\
\textbf{Key citations} & Vallerand 2003, Pelletier 1995, Fraser-Thomas 2006,
                          Wolpert 1992, Breiman 2001, Lundberg 2017 (SHAP) \\
\textbf{GitHub} & \url{https://github.com/KusalPabasara/sports-recommendation} \\
\end{tabular}
\end{keybox}

% ============================================================
\section{The Story to Tell --- One Narrative From Start to Finish}
% ============================================================

\begin{mdframed}[backgroundcolor=boxblue, linecolor=headerblue, linewidth=1.2pt,
                 innertopmargin=10pt, innerbottommargin=10pt,
                 innerleftmargin=12pt, innerrightmargin=12pt]
\textit{``The problem I started with is simple: 1.4 billion people are not
active enough, and part of the reason is that people struggle to find sports
they genuinely enjoy. The existing technology --- systems like AI Sport Scout
--- only look at physical fitness. But 30 years of sports psychology research
show that interest and motivation are the real drivers of whether someone
sticks with a sport.

So I built a machine learning system that flips this. Instead of asking how
fast you run, it asks what you enjoy --- team or individual? Indoor or
outdoor? High-energy or strategic? These interest answers, combined with
physical traits as secondary information, feed into a stacking ensemble of
XGBoost and Random Forest models. The ensemble outperforms every baseline
by 17.4\% on ranking quality.

But recommending sports you already know is only half the problem. If you
have only ever played cricket and football, the system might just recommend
cricket and football again. So I built a discovery score --- using cosine
similarity between your interest profile and the average interests of people
who play sports you have never tried. This achieves an 85.7\% novelty rate:
almost 9 out of every 10 recommendations are sports the user has never
encountered before.

To validate that the interest-first design is correct, I used SHAP to look
inside the model. It showed that interest and self-strength features account
for 65.9\% of the model's predictive power --- more than twice the
contribution of objective physical metrics. The model itself learned that
interests matter more.

The paper was accepted at ICDSIAI-26. The next step is a real-user study to
validate these findings with actual people, not synthetic data. This is
the proof-of-concept that makes that next step scientifically justified.''}
\end{mdframed}

\vspace{1cm}
\begin{center}
\small\color{gray}
This document was prepared as a full defence guide for the accepted paper.\\
All numbers are cross-verified against the submitted manuscript.\\
Paper: \texttt{paper-icdsiai/main.pdf} \quad|\quad
Code: \url{https://github.com/KusalPabasara/sports-recommendation}
\end{center}

\end{document}
